\section{The original arithmetic classes under Hilbert completion}
\label{CFStage:sec:fixed-packet-hilbert-geometry}

Fix the original complete counterfactual quartet $h$ of (HT1), the
integer $k\geq3$, and its complete cyclic annihilator
$\chi=\chi_{h,k}$ of degree $q\geq1$.
Every parameter remains fixed as $N\to\infty$. Retain
\[
 \begin{gathered}
 c=k/2,\quad \mathcal P=\mathbb C[S],\quad
 E=\mathcal P/\chi\mathcal P,\quad J:\mathcal P\longrightarrow E,\\
 \mathscr H_{h,k}=L^2(\mathbb R,m_{h,k}(u)\,du),\qquad
 \iota:\mathcal P\longrightarrow\mathscr H_{h,k},\quad
             (\iota P)(u)=P(c+iu).
 \end{gathered}
 \tag{FKG.1}
\]
The map $\iota$ is injective, since the original density is positive
almost everywhere and a nonzero polynomial has finitely many zeros
on the source line. We write $\mathcal P$ and $\chi\mathcal P$ for
these actual embedded subspaces when taking Hilbert-space quotients.
All quotient arrows below are complex linear. Multiplication by $S$
on the full Hilbert space has its unbounded-operator domain; no
everywhere-defined polynomial-module structure on that space is assumed.

\subsection{Vanishing of the entire quotient Gram}

Put $w_i=\lambda_i-c$ and $\delta_i=\Re w_i$ as in (FPK.14).
For a noncritical block, write $a_i(w)=C_{\epsilon_i}(w)^{-1}$,
where $\epsilon_i=\operatorname{sgn}\delta_i$. In normalized Taylor
coordinates the complete triangular matrix in (FPK.14) has entries
\[
 (\mathsf T_{N,i})_{j,l}=
 \begin{cases}
 N^{-\epsilon_iw_i/2}
 \displaystyle\sum_{r=0}^{j-l}
       a_i^{[r]}(w_i)
       \frac{(-\epsilon_i\log N/2)^{j-l-r}}{(j-l-r)!},&l\leq j,\\
 0,&l>j.
 \end{cases}
 \tag{FKG.2}
\]
This follows by the product rule for Taylor coefficients of
$a_i(w)\exp(-\epsilon_i w\log N/2)$; thus no derivative of the
original gamma unit has been dropped. For fixed roots and multiplicities,
\[
 \|\mathsf T_{N,i}\|
 \leq C_i N^{-|\delta_i|/2}(1+\log N)^{m_i-1}
       \longrightarrow0\qquad(\delta_i\ne0).
 \tag{FKG.3}
\]
For a critical block, (FPK.17) gives
$\|\mathsf T_{N,i}\|=(\log N)^{-1/2}\to0$ for $N\geq3$.
There are finitely many blocks, so $\|\mathsf T_N\|\to0$.

Let $\mathsf A_N=\mathsf T_N\mathsf K_N\mathsf T_N^*$.
The complete matrix limit (FPK.20) gives
$\mathsf A_N\to\mathsf H>0$. If $h_0$ is the smallest eigenvalue
of $\mathsf H$, then $\mathsf A_N\geq(h_0/2)I$ for all sufficiently
large $N$: for each unit vector the difference of the quadratic forms
is at most $\|\mathsf A_N-\mathsf H\|$. In particular
$\|\mathsf A_N^{-1}\|\leq2/h_0$. Undoing the exact congruence gives
\[
 \begin{aligned}
 \mathsf K_N^{-1}&=\mathsf T_N^*\mathsf A_N^{-1}\mathsf T_N,\\
 G_N^\sigma&=\mathsf E_\chi^*\mathsf T_N^*
                      \mathsf A_N^{-1}\mathsf T_N\mathsf E_\chi,\\
 \|G_N^\sigma\|&\leq\frac{2\|\mathsf E_\chi\|^2}{h_0}
                         \|\mathsf T_N\|^2\longrightarrow0.
 \end{aligned}
 \tag{FKG.4}
\]
All norms on $E$ here use the original increasing-power remainder
frame. In particular this is vanishing of the entire positive operator,
not merely of its determinant. The explicit bounds (FKG.3) also give
\[
 \|G_N^\sigma\|\leq C_\chi\left(
 \frac{\mathbf1_{\{\exists i:\delta_i=0\}}}{\log N}
 +\sum_{\delta_i\ne0}N^{-|\delta_i|}
                      (1+\log N)^{2m_i-2}\right)
 \tag{FKG.5}
\]
for all sufficiently large $N$, with one finite constant for the
fixed original observation.

Write $R_N^{\mathrm{ar}}:E\to\mathcal P_N$ for the canonical
least-norm polynomial lift in the original arithmetic form (AT6).
It is the unique minimizer on each affine fibre $J_NP=x$:
orthogonal projection onto $\chi\mathcal P_{N-q}$ in the finite
positive-definite polynomial space proves existence and uniqueness.
Its source Gram is $G_N^{\mathrm{ar}}$, so
\[
 J_NR_N^{\mathrm{ar}}=I_E,\qquad
 \|\iota R_N^{\mathrm{ar}}x\|_{\mathscr H_{h,k}}^2
                     =x^*G_N^{\mathrm{ar}}x.
 \tag{FKG.6}
\]
The proved actual envelope $m_{h,k}\leq A_k\sigma$ in (LC.24)
applies to every polynomial in the same affine fibre. Taking the
minimum first for the arithmetic norm and then comparing with the
minimizing gamma lift proves
\[
 0<G_N^{\mathrm{ar}}\leq A_kG_N^\sigma,\qquad
 \|G_N^{\mathrm{ar}}\|\to0,\qquad
 \|\iota R_N^{\mathrm{ar}}\|_{E\to\mathscr H_{h,k}}\to0.
 \tag{FKG.7}
\]
For the middle assertion a positive matrix has operator norm equal
to the supremum of its unit-vector quadratic forms; the last assertion
follows from (FKG.6). The constant $A_k$ is the original arithmetic
upper-comparison constant, including its full scalar mass.

\subsection{Exact arithmetic observations and theta primitives}

The original cohomology observation and its injection remain
\[
 \begin{gathered}
 \mathcal O=q^{(k)}\mathcal V:\mathcal P\longrightarrow Q^{\otimes k},
 \qquad \mathcal I=\sigma_h^{\otimes k}\eta_{h,k}:E\hookrightarrow Q^{\otimes k},\\
 \mathcal VP=P(D_1+\cdots+D_k)F_h^{\otimes k},\qquad
 \eta_{h,k}[P]=\upsilon_h^{\otimes k}P(A_h^{(k)})1,\\
 \upsilon_h=j_h(g/h),\quad g=2\xi,\quad D=-x\partial_x,
 \qquad\mathcal O=\mathcal I J.
 \end{gathered}
 \tag{FKG.8}
\]
These are precisely (AT2)--(AT2a), with the complete invertible
Taylor unit $\upsilon_h$ retained. Consequently for every $x\ne0$
and every $N\geq q-1$,
\[
 \mathcal O(R_N^{\mathrm{ar}}x)=\mathcal I x\ne0,
 \qquad \|\iota R_N^{\mathrm{ar}}x\|_{\mathscr H_{h,k}}\longrightarrow0.
 \tag{FKG.9}
\]
This comparison also retains the exact original tensor-source norm:
Mellin Plancherel on each variable gives
\[
 \|\mathcal VP\|_{L^2((0,\infty)^k,\,dx_1\cdots dx_k)}^2
 =\int_{\mathbb R^k}|P(k/2+i\textstyle\sum t_i)|^2
                         \prod_iw_h(t_i)\,d\mathbf t
 =\|\iota P\|_{\mathscr H_{h,k}}^2.
 \tag{FKG.10}
\]
Indeed the Mellin transform of the source is
$P(\sum s_i)\prod_iv_h(s_i)$, the $k$ Plancherel factors
$(2\pi)^{-k}$ are exactly those in $\prod_iw_h(t_i)$, and
the change $(t_1,\ldots,t_k)\mapsto(t_1,\ldots,t_{k-1},u)$
has determinant one. The final equality is the defining convolution.
Thus (FKG.9) quantifies the original source representatives themselves.

For $q-1\leq a<b$, monic division gives a unique polynomial
$\kappa_x\in\mathcal P_{b-q}$ satisfying
\[
 R_a^{\mathrm{ar}}x-R_b^{\mathrm{ar}}x=\chi\kappa_x.
 \tag{FKG.11}
\]
The inclusions into $\mathcal P_b$ are the literal coefficient
inclusions. Both remainders are $x$, so their difference is divisible
by $\chi$; the degree bound and uniqueness follow from monicity.
Division also proves linearity in $x$. To retain the entire original
boundary, successive monic division in $s_1,\ldots,s_k$ gives specific
polynomials $Q_i$ with
\[
 \chi(\textstyle\sum s_i)=\sum_{i=1}^kh(s_i)Q_i(\mathbf s).
 \tag{FKG.12}
\]
The final multivariable remainder is zero because the monomials of
degree less than $\deg h$ in each variable form a basis of
$(\mathbb C[s]/h)^{\otimes k}$ and the original $\chi$ annihilates
the sum operator on this algebra. Using the original
$h(D)F_h=\Theta\phi_0$, where
$\phi_0=(4\pi^2x^4-6\pi x^2)e^{-\pi x^2}$, form the cochain
\[
 \mathcal P_\chi(\kappa_x)=
 \sum_{i=1}^k(-1)^{i-1}Q_i(\mathbf D)
       \kappa_x(\textstyle\sum_jD_j)
 \bigl(F_h^{\otimes(i-1)}\otimes\phi_0
                          \otimes F_h^{\otimes(k-i)}\bigr).
 \tag{FKG.13}
\]
It has degree $k-1$ and
\[
 d_{\mathrm{tot}}\mathcal P_\chi(\kappa_x)
        =\mathcal V(R_a^{\mathrm{ar}}x-R_b^{\mathrm{ar}}x).
 \tag{FKG.14}
\]
At position $i$, the $i-1$ preceding degree-one factors contribute
$(-1)^{i-1}$ to the tensor differential, cancelling the displayed
sign. The relation $\Theta\phi_0=h(D_i)F_h$ then makes the sum
equal to (FKG.12) applied to the original tensor source. Polynomial
Euler operators commute with this differential by the original
identity $D\Theta=\Theta D$. This proves (FKG.14) term by term.
It is the full primitive (EW.47)--(EW.48), with the actual canonical
arithmetic lifts substituted.

\subsection{The exact algebraic quotient and its Hausdorff reduction}

Polynomials are dense in $\mathscr H_{h,k}$. Here the exponential
moment needed for the complete density theorem (ND.11)--(ND.12)
follows in the original sum coordinate: for $0<\eta<\pi/2$,
Tonelli's theorem and $|\sum t_i|\leq\sum|t_i|$ give
\[
 \int_{\mathbb R}e^{\eta|u|}m_{h,k}(u)\,du
 \leq\left(\int_{\mathbb R}e^{\eta|t|}w_h(t)\,dt\right)^k<\infty.
 \tag{FKG.15}
\]
Finiteness follows directly from (ND.3), choosing $\eta<2b<\pi/2$.
To specify the density argument in this one-dimensional application,
if $f$ is orthogonal to all polynomials, the finite complex measure
$d\nu=\overline f\,m_{h,k}\,du$ has an exponential moment of any
order less than $\eta/2$, by Cauchy--Schwarz and (FKG.15). Its
Fourier transform is holomorphic on the corresponding strip and
all its derivatives at zero vanish, because polynomials in $u$
are exactly polynomials in $S=c+iu$ under this invertible affine
coordinate map. Its Taylor series and then the identity theorem
make the transform zero on the strip. The Gaussian convolution
and approximate-identity proof (ND.12) then gives $\nu=0$ and
$f=0$ almost everywhere. Taking orthogonal complements proves
the asserted density, including the original measure and mass.

For any $P\in\mathcal P$ and $N\geq\max(q-1,\deg P)$, put
\[
 B_N(P)=P-R_N^{\mathrm{ar}}JP\in\chi\mathcal P.
 \qquad
 \|\iota(P-B_N(P))\|_{\mathscr H_{h,k}}\longrightarrow0.
 \tag{FKG.16}
\]
The first assertion follows by taking its remainder; the second is
(FKG.7) at the fixed vector $JP$. Thus the closure of
$\chi\mathcal P$ contains $\mathcal P$, whose density was just
proved. Consequently
\[
 \overline{\chi\mathcal P}^{\,\mathscr H_{h,k}}
       =\mathscr H_{h,k}.
 \tag{FKG.17}
\]
Every approximating relation in (FKG.16) has a literal polynomial
quotient by $\chi$, and hence the primitive (FKG.13) as well.
The isometry (FKG.10) also extends to the exact original tensor-source
completion:
\[
 \begin{gathered}
 \overline{\mathcal V}:\mathscr H_{h,k}\xrightarrow{\ \cong\ }
 \overline{\mathcal V(\mathcal P)}^{\,L^2((0,\infty)^k)},\\
 \overline{\mathcal V(\chi\mathcal P)}^{\,L^2((0,\infty)^k)}
       =\overline{\mathcal V(\mathcal P)}^{\,L^2((0,\infty)^k)}.
 \end{gathered}
 \tag{FKG.17a}
\]
Indeed polynomial approximants form Cauchy sequences in the original
tensor norm by (FKG.10); completeness defines the extension, and
that norm equality proves its independence of the approximants and
its isometry. Its range is closed by completeness and contains the
dense polynomial source, hence is the displayed closure. Applying
this isometry to (FKG.17) proves the second equality. Every vector
before closure on its left is the boundary (FKG.14) for its literal
polynomial quotient. No full test-space cohomology is quotiented
in forming these two specified Hilbert closures.

There is nevertheless an injective map of the unchanged algebraic
quotient into the algebraic Hilbert-space quotient, in the exact
sequence
\[
 \begin{gathered}
 0\longrightarrow E\xrightarrow{\ j\ }
       \mathscr H_{h,k}/\chi\mathcal P
       \xrightarrow{\ \pi\ }\mathscr H_{h,k}/\mathcal P
       \longrightarrow0,\\
 j([P])=\iota P+\chi\mathcal P,\qquad
 \pi(f+\chi\mathcal P)=f+\mathcal P.
 \end{gathered}
 \tag{FKG.18}
\]
The first map is independent of the polynomial representative.
Its kernel is zero because $\iota$ is injective and membership
in the embedded $\chi\mathcal P$ means literal polynomial
divisibility. The second map is well-defined and surjective since
$\chi\mathcal P\subset\mathcal P$; its kernel consists exactly of
the cosets represented by polynomials, which is the image of $j$.
This proves exactness at every term without a topological closure.

With the quotient seminorm inherited from the original Hilbert norm,
the Hausdorff reduction is instead the specific map
\[
 \begin{gathered}
 \mathscr H_{h,k}/\chi\mathcal P\longrightarrow
 \mathscr H_{h,k}/\overline{\chi\mathcal P}=0,\\
 f+\chi\mathcal P\longmapsto f+\overline{\chi\mathcal P}.
 \end{gathered}
 \tag{FKG.19}
\]
Indeed the seminorm of any coset is its distance to the dense
subspace $\chi\mathcal P$, hence zero. In particular (FKG.19)
kills the injected algebraic copy $j(E)$, although every nonzero
$x\in E$ retains the nonzero arithmetic class $\mathcal I x$ in
(FKG.8). This is an exact comparison between these two quotients.

One can also compute the entire closure of the original remainder
graph. Give $E$ its finite-dimensional topology in its original frame.
For any $f\in\mathscr H_{h,k}$ and $x\in E$, choose polynomials
$P_j\to f$ by density. For each fixed $j$, (FKG.7) permits an
$N_j\geq\max(q-1,\deg P_j)$ such that
$\|\iota R_{N_j}^{\mathrm{ar}}(x-JP_j)\|<1/j$. Then
\[
 \widetilde P_j=P_j+R_{N_j}^{\mathrm{ar}}(x-JP_j),\qquad
 J\widetilde P_j=x,\qquad\iota\widetilde P_j\longrightarrow f.
 \tag{FKG.20}
\]
It follows that the graph of $J:\mathcal P\subset\mathscr H_{h,k}\to E$
has closure $\mathscr H_{h,k}\times E$. In particular this densely
defined observation is not closable when $q\geq1$: (FKG.9) gives
sequences tending to zero with any prescribed nonzero remainder.
This precisely states why a continuous extension of the same
nonzero finite observation to this Hilbert completion is impossible.

At an original nonempty support label $\mathfrak a$, each linear
comparison just written lifts as
$(\mathfrak a,v)\mapsto(\mathfrak a,Lv)$ and sends the external
point $\tau$ to $\tau$. Thus the killed amplitude in (FKG.19) maps
to $(\mathfrak a,0)$, and the exact boundary (FKG.14) has that
same supported zero observation. The label remains specified.
These maps quantify this particular Hilbert completion; they do
not identify its zero Hausdorff quotient with the entire original
arithmetic cohomology or alter any original support restriction.

\subsection{The coefficient for the complete counterfactual quartet}

In the actual off-line counterfactual (HT1)--(HT2), retain
$\rho=1/2+\delta+i\gamma$ with $0<\delta<1/2$, $\gamma>0$,
and its full order $m$. Put $\ell_k=1+k(m-1)$. The distinct
cyclic roots and their complete multiplicities are exactly
\[
 \lambda_{a,b}=k/2+(2a-k)\delta+i(2b-k)\gamma,
 \quad 0\leq a,b\leq k,\qquad m_{a,b}=\ell_k.
 \tag{FKG.21}
\]
For completeness, in each ordered tensor-local factor the nilpotent
sum $n_1+\cdots+n_k$, with $n_i^m=0$, has nonzero top coefficient
$(k(m-1))!/((m-1)!)^k$ at total degree $k(m-1)$ and its next
power vanishes. Its index is therefore $\ell_k$. The real and
imaginary signs may be chosen independently in each tensor factor,
so every pair $(a,b)$ occurs. Since $\delta\gamma\ne0$, different
pairs give different sums. The full annihilator is consequently
the product over (FKG.21), each to power $\ell_k$, with
$q=\ell_k(k+1)^2$ as originally specified.

The exponents of the full gamma determinant therefore evaluate to
\[
 \begin{aligned}
 \mathcal A_\chi
 &=\delta\ell_k(k+1)\sum_{a=0}^k|2a-k|
   =2\delta\ell_k(k+1)
              \left\lfloor\frac{(k+1)^2}{4}\right\rfloor
   =\mathcal L_{h,k},\\
 \mathcal B_\chi
 &=\begin{cases}(k+1)\ell_k^2,&k\text{ even},\\0,&k\text{ odd}.
   \end{cases}
 \end{aligned}
 \tag{FKG.22}
\]
To verify the finite sum, pair $a$ with $k-a$. For $k=2r$ the
sum of its positive terms is $2(1+\cdots+r)=r(r+1)$; for
$k=2r-1$ it is $1+3+\cdots+(2r-1)=r^2$. Doubling gives
$2\lfloor(k+1)^2/4\rfloor$ in both cases. A root is critical
exactly when $2a=k$, which occurs only for even $k$, with its
$k+1$ distinct imaginary coordinates and multiplicity $\ell_k$.
This proves the second exponent with all nilpotent jets included.
The equality with $\mathcal L_{h,k}$ uses its exact definition
in (HT8), not a change of normalization.

In particular (LC.26) gives, on the actual unchanged arithmetic
source for every fixed $k\geq3$,
\[
 \lim_{N\to\infty}\frac{\log V_N^{\mathrm{ar}}}{\log N}
 =-2\delta\ell_k(k+1)
              \left\lfloor\frac{(k+1)^2}{4}\right\rfloor.
 \tag{FKG.23}
\]
The value of $\mathcal B_\chi$ in (FKG.22) is the proved gamma
$\log\log N$ coefficient. For a mixed critical and noncritical
arithmetic packet, the envelope proof establishes (FKG.23), and
does not establish its next $\log\log N$ coefficient. None of
(FKG.3)--(FKG.23) supplies an error uniform in growing $k$, $q$,
or the packet. Their precise conclusions concern fixed original
objects and the displayed exact maps.
